% Computer Science A --- testing, mocks, logging, and SRE/JUnit literacy (not Spring-specific).
% Free-response block migrated from \texttt{cs-a/04-practice.tex}; multiple-choice block from \texttt{cs-a/10-practice.tex}.
% Included from \texttt{main.tex} in the CS-A block (after \texttt{functional}, before \texttt{cs-db}).

\runningheader{Computer Science A}{Testing, mocks, and logging}{Page \thepage\ of \numpages}

\begingradingrange{csau4}

\uplevel{\par\textbf{Testing and mocks}\par}

\question[4]
What is Mockito and why would we use it?
\begin{solutionorbox}[1.25in]
  Popular Java mocking library: stub dependencies, verify interactions, isolate the class under test without slow or brittle real I/O or databases.
\end{solutionorbox}

\question[4]
What is a mock?
\begin{solutionorbox}[1.05in]
  A test double with programmed behavior and expectations: returns canned values, records calls, and fails tests if expected interactions do not occur.
\end{solutionorbox}

\question[4]
What is test driven development?
\begin{solutionorbox}[1.2in]
  Write a failing test first, implement minimal code to pass, refactor; repeats in short cycles to drive design and regression safety.
\end{solutionorbox}

\question[4]
Why are unit tests important?
\begin{solutionorbox}[1.25in]
  Catch regressions early, document expected behavior, enable safe refactoring, shorten debug loops, and support continuous integration.
\end{solutionorbox}

\question[4]
How can JUnit annotations help with running our tests?
\begin{solutionorbox}[1.35in]
  Mark lifecycle (\texttt{@BeforeEach}, \texttt{@AfterEach}), identify tests (\texttt{@Test}), disable/skip, parameterize, order, tag for suites, and set timeouts or expected exceptions.
\end{solutionorbox}

\question[4]
Give me an example of how you would test a method that takes in two integers as input and returns an integer representing the sum.
\begin{solutionorbox}[1.55in]
  Call the method with known pairs and assert the result, e.g.\ \texttt{assertEquals(5, add(2, 3));} plus edge cases (\texttt{0}, negatives) using JUnit \texttt{assertions}.
\end{solutionorbox}

\newpage

\uplevel{\par\textbf{Logging}\par}

\question[4]
What is logging and what are the benefits of it?
\begin{solutionorbox}[1.25in]
  Recording runtime events (errors, timing, decisions) to appenders (console, files, aggregators). Benefits: production diagnosis without debugger, audit trails, and structured observability.
\end{solutionorbox}

\question[4]
Describe the different logging levels and how they should be used.
\begin{solutionorbox}[1.45in]
  Typical levels (e.g.\ SLF4J): \texttt{ERROR}/\texttt{WARN} for problems, \texttt{INFO} for high-signal lifecycle events, \texttt{DEBUG}/\texttt{TRACE} for developer detail---tune thresholds per environment to limit noise and cost.
\end{solutionorbox}

\question[4]
How would we configure logging in an application?
\begin{solutionorbox}[1.35in]
  Use a facade (\texttt{slf4j}) with a binding (Logback, Log4j2): \texttt{logback.xml} or \texttt{log4j2.xml} for appenders, levels per logger, patterns, and optional env-specific profiles.
\end{solutionorbox}

\endgradingrange{csau4}

\newpage

\begingradingrange{csau10}

\question[4]
Henry came up with a few options for logging and presented it to Eric the SRE team lead. Which of the following option would be of highest priority to log?
\begin{choices}
  \choice 200 status codes
  \choice When the application is on low usage
  \CorrectChoice When the end user experiences a slow-down of the application
  \choice None of the above
\end{choices}
\begin{solution}
  User-visible latency and errors usually trump routine success metrics for incident response.
\end{solution}

\question[4]
When you write a test case, you should test every single method of an application.
\begin{choices}
  \CorrectChoice False
  \choice True
\end{choices}
\begin{solution}
  Prefer focused tests: one behavior or rule per test; use several \texttt{@Test} methods rather than one mega-test per production method.
\end{solution}

\question[4]
\texttt{@BeforeEach} will run every time, before each unit test.
\begin{choices}
  \CorrectChoice True
  \choice False
\end{choices}
\begin{solution}
  JUnit~5 \texttt{@BeforeEach} (JUnit~4 \texttt{@Before}) runs before each test method in that class.
\end{solution}

\question[4]
What is the package for JUnit?
\begin{choices}
  \choice \texttt{java.junit}
  \CorrectChoice \texttt{org.junit}
  \choice \texttt{javax.junit}
  \choice \texttt{org.oracle.java.junit}
  \choice \texttt{java.test.unit}
\end{choices}
\begin{solution}
  JUnit~4 lives under \texttt{org.junit}; JUnit~5 APIs are primarily \texttt{org.junit.jupiter}.
\end{solution}

\newpage

\question[4]
Predict the output of the following test:
\begin{lstlisting}[style=java]
public class ExpenseApprover {
  public boolean approveOrDenyExpense(Expense e) {
    if (e.amount < 100 && e.hasReceipt) {
      return true;
    } else {
      return false;
    }
  }
}

public class Test {
  ExpenseApprover ea = new ExpenseApprover();

  @Test
  public void testExpense() {
    Expense exp = new Expense();
    exp.setAmount(105);
    exp.setReceipt(true);
    Assert.assertFalse(ea.approveOrDenyExpense(exp));
  }
}
\end{lstlisting}
\begin{choices}
  \CorrectChoice Pass
  \choice Fail
  \choice Compilation error
  \choice Test doesn't run
\end{choices}
\begin{solution}
  Amount 105 fails \texttt{amount < 100}, so the method returns \texttt{false}; assertion expects \texttt{false}---passes.
\end{solution}

\endgradingrange{csau10}
